\begin{abstract}
Industry systems already train text generators against predictions of real clicks. Whether this makes headlines less true has been tested only with simulated audiences or in short ablations on private data. We test it on public data. We use reinforcement learning (RL) to train a 3B headline writer for news articles under three rewards: a click predictor fitted to randomised headline tests from the Upworthy archive, a simulated audience of LLM personas, and the click predictor with a faithfulness penalty. The study was pre-registered with two seeds per reward, and every headline is checked against its article. Under both the click predictor and the simulated audience, headlines came to claim things their article does not say: invented reactions, roles and events. Only the click predictor, which also moved the writer four to six times further from its starting point, went further. The writer exploited it by repeating one fabricated viral story, whatever the article: invented-number flags rose from 1.7\% to 98\%, against about 2\% under the other rewards. Best-of-16 selection with the same predictor did not produce this story. The faithfulness penalty reduced invented claims but did not remove them. An LLM probe missed almost every fabrication, and in labelled samples from training, number and entity checks missed most invented events and roles. Fabrication checks for engagement-trained writers must compare each claim with the article.
\end{abstract}
